\documentclass[11pt,a4paper]{article}

\usepackage{thga-db}

\renewcommand{\thgaKurs}{Databases}
\renewcommand{\thgaDatum}{Summer Term 2026}

\begin{document}

\thispagestyle{empty}
\thgaTitel{Term Project -- Documentation}%
{Chess Club Tournament \& Rating Manager}

\begin{infobox}
\textbf{Author:} Yahya El Idrissi \\
\textbf{Module:} Introduction to Database Management Systems -- THGA Bochum \\
\textbf{Stack:} PostgreSQL 16 -- FastAPI -- Docker Compose -- tkinter (\texttt{.deb}) \\
\textbf{Application repository:} \url{https://github.com/yahya2912/chess-club-manager} \\
\textbf{Documentation repository:} \url{https://github.com/yahya2912/chess-tournament-documentation} \\
\textbf{Demo video:} \url{https://youtu.be/xtMlvPf5F0g}
\end{infobox}

\tableofcontents

\vspace{1em}

% ============================================================
\section{Introduction and motivation}
% ============================================================

Chess clubs often organize small tournaments with spreadsheets, handwritten
pairing sheets, or several independent tools. This makes it difficult to keep
player data, tournament registrations, game results, Elo changes, and final
standings consistent.

The \textbf{Chess Club Tournament \& Rating Manager} replaces this fragmented
workflow with a relational database-backed desktop application. It manages
\textbf{players}, \textbf{tournaments}, \textbf{registrations},
\textbf{rounds}, \textbf{games}, \textbf{openings}, and
\textbf{rating history}. Data is exposed through a FastAPI REST API and used by
a tkinter desktop frontend.

\begin{beispielbox}
\textbf{Usage scenario.}
A club administrator creates a tournament, registers the participating players,
records the result of a game, and immediately sees the updated Elo ratings and
tournament standings. The rating changes are also written to the rating-history
table so that previous values remain traceable.
\end{beispielbox}

The delivered core workflow is:

\begin{verbatim}
Create players
-> Create tournament and rounds
-> Register players
-> Record games
-> Update Elo and rating history
-> View standings
\end{verbatim}

Automatic pairing generation is intentionally outside the project scope.

% ============================================================
\section{Data model}
% ============================================================

The domain contains seven entity types:

\begin{itemize}[leftmargin=1.4em, topsep=2pt]
  \item \textbf{player}
  \item \textbf{opening}
  \item \textbf{tournament}
  \item \textbf{registration}
  \item \textbf{round}
  \item \textbf{game}
  \item \textbf{rating\_history}
\end{itemize}

A player can participate in many tournaments, and a tournament can contain many
players. This \texttt{N:M} relationship is resolved by
\textbf{registration}. A tournament contains several rounds, and each game
belongs to one round. Every game references a White player and a Black player
and may also reference an opening through its ECO code. Rating changes are
stored separately in \textbf{rating\_history} so that previous values are not
overwritten.

\begin{beispielbox}
\textbf{ER model (textual).}

\begin{itemize}[leftmargin=1.4em, topsep=2pt]
  \item \textbf{player} (1) -- registers -- (N) \textbf{registration}
  \item \textbf{tournament} (1) -- contains -- (N) \textbf{registration}
  \item \textbf{tournament} (1) -- contains -- (N) \textbf{round}
  \item \textbf{round} (1) -- contains -- (N) \textbf{game}
  \item \textbf{player} (1) -- plays as White/Black -- (N) \textbf{game}
  \item \textbf{player} (1) -- has -- (N) \textbf{rating\_history}
  \item \textbf{opening} (1) -- classifies -- (N) \textbf{game}
\end{itemize}

Thus \textbf{player} \texttt{N:M} \textbf{tournament}, resolved through
\textbf{registration}.
\end{beispielbox}

\subsection{Relational schema}

The implemented schema in \texttt{db/01\_schema.sql} uses the following
relations:

\begin{itemize}[leftmargin=1.4em]
  \item \texttt{player}(\underline{id}, name, birth\_year, elo)
  \item \texttt{opening}(\underline{eco\_code}, name)
  \item \texttt{tournament}(\underline{id}, name, start\_date, status, rounds)
  \item \texttt{registration}(\underline{player\_id, tournament\_id}, seed)
  \item \texttt{round}(\underline{id}, tournament\_id, round\_number)
  \item \texttt{game}(\underline{id}, round\_id, white\_id, black\_id,
        opening\_eco, result)
  \item \texttt{rating\_history}(\underline{id}, player\_id, game\_id,
        old\_elo, new\_elo, changed\_at)
\end{itemize}

Important integrity rules include primary and foreign keys, a composite key for
registrations, unique round numbers within a tournament, controlled tournament
status values, valid game results, and a check that White and Black are not the
same player.

\textbf{Normalisation.}
Player data is stored only in \texttt{player}, opening descriptions only in
\texttt{opening}, tournament data only in \texttt{tournament}, and game,
registration, round, and rating-history data reference those entities through
keys instead of duplicating their attributes. The design therefore avoids
unnecessary redundancy and keeps the main relations in third normal form.

% ============================================================
\section{Database (PostgreSQL)}
% ============================================================

The schema and deterministic seed data are stored in:

\begin{verbatim}
db/01_schema.sql
db/02_seed.sql
\end{verbatim}

PostgreSQL is started through Docker Compose. The schema is initialized when
the database is created, and the seed data provides deterministic reference and
sample rows that are also used by CI checks.

The database is not treated as a passive store. Integrity rules are enforced
directly at DDL level so invalid states are rejected even when they do not
originate from the graphical frontend.

Examples include:

\begin{itemize}[leftmargin=1.4em]
  \item one registration per player and tournament,
  \item unique round numbers inside a tournament,
  \item restricted game-result values,
  \item White and Black must be different players,
  \item foreign keys between tournaments, rounds, games, players, openings,
        registrations, and rating history.
\end{itemize}

Multi-row operations are executed transactionally in the backend. This is
especially important when recording a game because one user action inserts the
game, updates two current Elo values, and inserts two rating-history entries.

% ============================================================
\section{REST API (FastAPI)}
% ============================================================

The FastAPI backend is the only application component that communicates with
PostgreSQL. The tkinter frontend never executes SQL directly.

The API entry point is:

\begin{verbatim}
api/app/main.py
\end{verbatim}

Application logic is separated into routers for players, tournaments,
registrations, games, rating history, and standings.

The public \texttt{/health} endpoint is used to verify that the backend is
reachable. Application routes are protected with the
\texttt{X-API-Key} header.

\renewcommand{\arraystretch}{1.2}
\begin{tabularx}{\linewidth}{@{}p{1.6cm}p{4.7cm}X@{}}
\toprule
\textbf{Method} & \textbf{Path / area} & \textbf{Purpose} \\
\midrule
\texttt{GET} & \texttt{/players} & list players and current Elo values \\
\texttt{POST} & \texttt{/players} & create a player and initial rating history \\
\texttt{POST} & \texttt{/tournaments} & create tournament and round records \\
\texttt{POST} & \texttt{/registrations} & register a player for a tournament \\
\texttt{POST} & \texttt{/games} & validate and store a game, then update Elo \\
\texttt{GET} & \texttt{/rating-history} & return Elo history, optionally filtered \\
\texttt{GET} & \texttt{/standings} & calculate tournament standings and Buchholz \\
\bottomrule
\end{tabularx}

\subsection{Elo calculation}

Elo calculation is isolated in:

\begin{verbatim}
api/app/elo.py
\end{verbatim}

The delivered implementation uses a K-factor of 20. After a valid game, both
players receive new ratings and two history rows are inserted.

\subsection{Standings and Buchholz}

Standings are calculated in SQL using common table expressions. Scoring is:

\begin{itemize}[leftmargin=1.4em]
  \item win = 1.0 point,
  \item draw = 0.5 points,
  \item loss = 0.0 points.
\end{itemize}

Full Buchholz is calculated as the sum of all opponents' tournament points.
The final ordering uses points first, followed by Buchholz and wins.

\begin{lstlisting}[language=SQL, caption={Simplified scoring logic used for standings}]
CASE g.result
    WHEN '1-0' THEN 1.0
    WHEN '0-1' THEN 0.0
    ELSE 0.5
END AS points
\end{lstlisting}

\subsection{Extending the backend}

To add a new function, the same separation should be preserved:

\begin{enumerate}[leftmargin=1.6em, topsep=2pt]
  \item decide whether the database schema must change,
  \item add or extend a FastAPI route,
  \item reuse the API-key dependency,
  \item use a transaction for multi-row writes,
  \item extend the frontend API client,
  \item add the graphical control or view,
  \item add tests for valid and invalid cases.
\end{enumerate}

% ============================================================
\section{Frontend (tkinter, packaged as .deb)}
% ============================================================

The frontend is a Python tkinter desktop application. It communicates only with
the REST API and does not connect directly to PostgreSQL.

At startup, a connection dialog asks for the API URL and the
\texttt{X-API-Key}. The connection is validated before the main window opens.

The main application contains six tabs:

\begin{itemize}[leftmargin=1.4em]
  \item \textbf{Players}
  \item \textbf{Standings}
  \item \textbf{Rating History}
  \item \textbf{Record Game}
  \item \textbf{Register Player}
  \item \textbf{Create Tournament}
\end{itemize}

\subsection{User workflow}

\textbf{Players.}
Existing players can be viewed and new players can be created with a name,
optional birth year, and initial Elo value. Creating a player also creates the
first rating-history entry.

\textbf{Create Tournament.}
The user enters a tournament name, start date, and number of rounds. The backend
creates the tournament and its rounds in one operation.

\textbf{Register Player.}
A player is assigned to a tournament with an optional seed. Duplicate
registration is rejected.

\textbf{Record Game.}
The user selects the round, White player, Black player, ECO code, and result.
Valid results are \texttt{1-0}, \texttt{0-1}, and
\texttt{1/2-1/2}. A successful request stores the game and updates both
ratings.

\textbf{Standings.}
The standings view displays points, Buchholz, wins, draws, losses, and games
played.

\textbf{Rating History.}
The rating-history view shows all recorded Elo changes or filters them by
player.

\subsection{Packaging}

For Debian-based systems the frontend is packaged as a \texttt{.deb}. The
package can be built with:

\begin{verbatim}
sh frontend/debian/build-deb.sh
\end{verbatim}

The generated package is:

\begin{verbatim}
frontend/debian/chess-club-manager_1.0.0_all.deb
\end{verbatim}

The package contains the desktop frontend only. PostgreSQL and FastAPI remain
containerized and are started separately through Docker Compose.

% ============================================================
\section{Operation and deployment}
% ============================================================

The backend runs as two containers, PostgreSQL and FastAPI, orchestrated by
Docker Compose. Credentials and the API key are stored in a local
\texttt{.env} file that is not committed.

\subsection{Local startup}

Create the local configuration:

\begin{verbatim}
cp .env.example .env
\end{verbatim}

Start the backend:

\begin{verbatim}
docker compose up -d --build
\end{verbatim}

Check the health endpoint:

\begin{verbatim}
http://127.0.0.1:8000/health
\end{verbatim}

A working backend returns:

\begin{verbatim}
{"status":"ok"}
\end{verbatim}

Start the source frontend on Linux or macOS:

\begin{verbatim}
PYTHONPATH="$PWD/frontend" python3 -m chess_club_fe.main
\end{verbatim}

At startup, enter the API URL and the \texttt{API\_KEY} value from
\texttt{.env}.

\subsection{Deployment steps}

On the target system:

\begin{enumerate}[leftmargin=1.6em, topsep=2pt]
  \item copy the project and create the \texttt{.env} file,
  \item run \texttt{docker compose up -d --build},
  \item install the generated Debian package with \texttt{dpkg},
  \item launch \texttt{chess-club-manager},
  \item enter the API URL and matching \texttt{X-API-Key}.
\end{enumerate}

\subsection{Testing and evaluation}

For local API tests:

\begin{verbatim}
python3 -m venv api/.venv
api/.venv/bin/pip install -r api/requirements.txt
api/.venv/bin/pip install -r api/requirements-dev.txt
cd api
.venv/bin/pytest tests -v
\end{verbatim}

GitHub Actions verifies the PostgreSQL schema and deterministic seed data,
database acceptance checks, FastAPI endpoint and business-logic tests, Docker
Compose startup, Debian package creation, and LaTeX documentation builds.

The main functions were verified as follows:

\begin{center}
\begin{tabular}{>{\raggedright\arraybackslash}p{3.0cm}
                >{\raggedright\arraybackslash}p{4.7cm}
                >{\raggedright\arraybackslash}p{5.0cm}}
\toprule
\textbf{Function} & \textbf{Test} & \textbf{Expected result} \\
\midrule
Create player &
Create a player with an initial Elo value &
Player and initial rating-history row are stored. \\

Create tournament &
Create a tournament with multiple rounds &
Tournament and requested round records are created. \\

Register player &
Register the same player twice &
First registration succeeds; duplicate registration is rejected. \\

Record game &
Record a valid result between registered players &
Game is stored, both Elo values change, and history rows are created. \\

Invalid game &
Use an invalid result or identical White/Black player &
Request is rejected and no partial update remains. \\

Standings &
Open standings after recorded games &
Points, Buchholz, wins, draws, losses, and games played are shown. \\

Authentication &
Call a protected route with a wrong API key &
Request is rejected. \\
\bottomrule
\end{tabular}
\end{center}

\subsection{Documentation build and release}

This documentation repository uses the professor's LaTeX workflow structure.
The local build command is:

\begin{verbatim}
make
\end{verbatim}

The resulting file is:

\begin{verbatim}
out/documentation.pdf
\end{verbatim}

A version tag matching \texttt{v*} triggers the GitHub Actions release
workflow. For example:

\begin{verbatim}
git tag v1.0.0
git push origin v1.0.0
\end{verbatim}

The generated documentation PDF is then attached to the GitHub Release.

% ============================================================
\section{Reflection}
% ============================================================

The separation into database, API, and frontend made the project easier to
develop and test. Database constraints can be checked independently, API logic
can be tested before the GUI is involved, and the tkinter frontend does not
contain SQL.

The most consistency-sensitive operation is game recording. One action affects
the game table, two current Elo values, and two rating-history entries. Using a
transaction prevents partial updates if validation or a database operation
fails.

The deployment model also influenced the design. PostgreSQL and FastAPI are
reproducible through Docker Compose, while the tkinter frontend remains outside
Docker and can be installed as the required Debian package.

The delivered version intentionally does not include automatic Swiss pairing,
edit/delete operations through the GUI, automatic tournament closing, user
accounts, or detailed opening analytics. These are reasonable future
extensions but were kept outside the defined scope.

% ============================================================
\section*{Sources and reuse}
% ============================================================

\addcontentsline{toc}{section}{Sources and reuse}

Course patterns were reused and adapted where appropriate:

\begin{itemize}[leftmargin=1.4em, topsep=2pt]
  \item PostgreSQL/FastAPI separation and Docker Compose backend structure --
        adapted from the DBMS lecture and exercises.
  \item \texttt{X-API-Key} protection -- adapted from the course API example.
  \item tkinter desktop-client approach and Debian packaging workflow --
        adapted from the course material.
  \item Documentation structure, THGA style, Makefile, and GitHub Actions
        release workflow -- adapted from the professor's DBMS\_10 template.
\end{itemize}

The chess-specific schema design, Elo calculation, tournament workflow,
standings and Buchholz logic, tests, and application-specific documentation
were developed for the Chess Club Tournament \& Rating Manager.

\end{document}
